\section{Introduction}
\label{sec:introduction}

Power-flow feasibility asks whether the electrical laws and operating limits
of a network can hold simultaneously. Even without switching decisions,
power is nonlinear in voltage. This leaves two separate questions: how hard
is it to decide exact feasibility, and how much does a small numerical
residual establish? We answer both for a resistive model with rational
conductances, positive voltage intervals, and signed injection intervals,
and transfer the exact classification to several precisely specified AC
models. The results identify worst-case obstructions to exact certification;
they make no claim about the typical difficulty of operating networks.

Our main result is that exact resistive feasibility is complete for
$\ER$, the class associated with the existential theory of the reals.
For each fixed integer $r$, hardness holds simultaneously on connected
simple planar bipartite graphs of maximum degree three and girth at least
$r$, with unit conductance and interval endpoints from fixed finite rational
sets depending only on $r$ (Theorem~\ref{thm:structural}). Thus neither
large numerical coefficients nor dense graphs are required. The class
satisfies $\NP\subseteq\ER\subseteq\PSPACE$; its relation to exact
real-valued certificates and its range of applications are surveyed by
\citet{SchaeferEtAl2024}. A polynomial-size certificate scheme for every
feasible instance, with deterministic polynomial-time verification, would
therefore imply $\NP=\ER$. The separation of these classes is open.

The principal contribution is an electrical realization of bounded arithmetic
that preserves the entire solution set, together with transformations that
impose all the graph and data restrictions above. To the best of our
knowledge, the prior power-flow results discussed below do not provide such
a realization or an $\ER$-completeness theorem for this restricted resistive
model. The realization also supports two further conclusions: restricted
networks can have arbitrary compact semialgebraic feasible-set topology, and
an infeasible network can admit doubly exponentially small injection
residuals. The AC results specify which angle conventions preserve these
conclusions. They use established cycle and winding ideas; the contribution
there is their exact polynomial encoding and the resulting complexity
transfers.

\subsection{Results, significance, and proof strategy}

The results have three connected parts.

\begin{enumerate}
\item \emph{Exact complexity and electrical arithmetic.}
A bus with prescribed voltage and injection imposes an affine relation on
neighboring voltages.
Alternating complement paths supply distinct copies of a source variable
without increasing the degree beyond three. A small nonlinear gadget
then enforces inversion. Every source solution has exactly one network
extension, and designated bus voltages recover the source variables
(Section~\ref{sec:resistive}). A bounded arithmetic crossover gives
planarity; free-injection connectors give connectedness; and harmonic
subdivisions give unit conductance, bipartiteness, and prescribed girth
(Section~\ref{sec:structural}). All these transformations preserve the
rational solution-set correspondence. This identifies a source of hardness
that survives simultaneous sparsity and coefficient restrictions. The hard
graphs may have many cycles: large girth does not mean that they are trees.

\item \emph{AC angle semantics and solution sets.}
A local cosine inequality constrains the principal angle of a line but
need not admit real bus angles whose ordinary differences satisfy the same
limit. A cycle can wind around the origin while every local inequality
holds. We encode the additional zero-winding requirement by a quadratic
existential formula with $O(|N|+|E|)$ variables, predicates, and monomial
occurrences, including unequal voltage magnitudes and every rational cosine
in $(-1,1]$ (Section~\ref{sec:ac}). For resistive
lines, an energy identity forces equal real angles when all reactive
injections vanish and all line differences have magnitude less than $\pi$.
This transfers the resistive classification to AC feasibility, including
one-sided reactive intervals of positive width. Separate transfers cover
reference-fixed bus-angle boxes and strictly positive principal-angle
windows that shrink with the number of buses. The latter use
polynomial-bit rational cosines, which need not come from a fixed finite set.

The preserved solution sets yield rational equivalence precisely for compact
basic closed sets over $\Q$, that is, sets described by a conjunction of
polynomial equalities and weak inequalities with rational coefficients
(Section~\ref{sec:algebraic}). Arbitrary compact semialgebraic sets still
occur up to semialgebraic homeomorphism. A rational network can also have a
unique feasible voltage vector with a designated voltage generating any
prescribed real algebraic number field. These results concern topology and
exact algebraic coordinates, beyond a count of isolated solutions.

\item \emph{Limits of residual-based certification.}
The ordinary electrical reduction satisfies a two-sided residual bound.
An explicit infeasible family with $O(k)$ buses, degree at most three, and
fixed numerical data admits injection residuals below $2^{-2^k}$ while
respecting every voltage bound exactly (Section~\ref{sec:numerical}).
A general polynomial-minimum bound gives the matching doubly exponential
scale for this fixed-data bounded-degree class. The residual family is not
asserted to have all the additional restrictions of the planar construction.
A universal residual threshold for exact feasibility can therefore require
exponentially many accuracy bits in the number of buses. This is a limitation
of certification by residual size alone, not a running-time lower bound for
all exact algorithms. At an explicit gap promise, short rational certificates
do exist. An energy estimate also bounds angle deviation and active-power
error caused by small reactive injections.
\end{enumerate}

\subsection{Prior power-flow results and model scope}

The resistive equations are the physical direct-current power-flow model
studied by \citet{GanLow2014}. They differ from the linear ``DC
approximation'' of AC power flow. Gan and Low give sufficient conditions for
exactness of a second-order cone relaxation of resistive optimal power flow,
including conditions involving voltage upper bounds and injection lower
bounds. Our classification allows independent voltage and injection
intervals, both of which may be singletons at the same bus. In particular,
the reduction uses buses with both prescribed voltage and prescribed
injection. It does not assume the hypotheses of those relaxation-exactness
results. Section~\ref{sec:foundations} states these input choices explicitly;
we do not establish hardness after removing the prescribed quantities used
by the construction.

A related resistive model fixes source voltages and asks which
constant-power demands can be supplied at load buses.
\citet[Theorems 3.18 and 3.22]{JeeningaEtAl2023} characterize exact feasibility and
its interior by matrix-inequality alternatives, and establish convexity of
the feasible demand set. Their model has no independent operating voltage
bounds or injection restrictions at the fixed-voltage sources. Those
restrictions distinguish our input model; signed demands alone do not.
An exact matrix-inequality characterization also does not, by itself,
establish a polynomial-time Turing algorithm for exact feasibility.

For AC networks, \citet{BienstockVerma2019} prove strong NP-hardness in a
lossless model with fixed voltage magnitudes and unconstrained reactive
power. \citet{LehmannEtAl2016} prove NP-hardness on trees, already using star
networks, in a fixed-magnitude model with active and reactive constraints.
Our classification identifies the full existential-real complexity of a
different subclass: resistive lines and variable magnitudes, with reactive
constraints that force equal angles under the stated angle convention.
It neither strengthens the tree restriction in their theorem nor proves
$\ER$-hardness for the lossless fixed-magnitude model. The model
restrictions are summarized in
Table~\ref{tab:prior-complexity}.

\begin{table}[t]
\centering
\small
\caption{Exact feasibility classifications for different power-flow models.
The restrictions within each row hold simultaneously.}
\label{tab:prior-complexity}
\begin{tabular}{@{}p{0.23\linewidth}p{0.49\linewidth}p{0.22\linewidth}@{}}
\toprule
Result & Model and restrictions & Classification \\
\midrule
\citet{BienstockVerma2019} & Lossless AC; fixed voltage magnitudes;
unconstrained reactive power & Strongly NP-hard \\
\addlinespace
\citet{LehmannEtAl2016} & AC on stars; fixed voltage magnitudes;
active and reactive constraints & NP-hard \\
\addlinespace
Theorem~\ref{thm:structural} & Resistive; connected, simple, planar,
bipartite; degree at most three; any fixed girth bound; unit conductance;
fixed finite endpoint sets & $\ER$-complete \\
\addlinespace
Theorem~\ref{thm:ac-complete} & Resistive-line AC; variable magnitudes;
zero reactive injections; real line differences in $[-\pi/2,\pi/2]$
& $\ER$-complete \\
\bottomrule
\end{tabular}
\end{table}

Cycle consistency and winding are established features of power-flow
models. \citet[Theorem 2]{FarivarLow2013} characterize branch-flow angle
recovery by cycle sums that vanish modulo $2\pi$.
\citet{DelabaysEtAl2016} relate loop currents to winding numbers, and
\citet{JafarpourEtAl2022} develop winding cells and uniqueness, modulo a
common rotation, within each cell for monotone flow networks on the torus.
In our real-angle model, the
chosen short line differences must instead sum to zero as real numbers.
This selects the zero-winding cell. We do not claim that criterion, the
existence of winding solutions, or uniqueness in the monotone angular
regime as new. Our additional step is to express the required branch
choices and zero winding in polynomial size with rational coefficients,
without taking an inverse trigonometric function as an exact arithmetic
primitive. This makes the $\ER$ membership statement precise.

The resistive-AC connection with zero reactive injections already appears in
\citet[Appendix B, Case 2]{LavaeiLow2012}; weighted sine-coupling equilibria
also arise in \citet{DorflerEtAl2013}. We give a self-contained energy proof
under the real-angle hypothesis needed here. On meshed networks, the
principal-angle winding examples in Section~\ref{subsec:principal} have
zero reactive injection and nonconstant phasors. Thus the angle convention
is material to the transfer, rather than a choice of coordinates.

\subsection{Algebraic and quantitative precedents}

The algebraic study of power flow predates this work.
\citet{MarecekEtAl2016} study polynomial formulations, generic complex
solution counts, and positive-dimensional solution sets. Our universality
results instead realize prescribed compact feasible sets, with operating
inequalities, in the restricted resistive family. They do not assert that
such behavior is generic. The arithmetic starting problem is due to
\citet{AbrahamsenEtAl2018,AbrahamsenEtAl2022}, and the planar crossover is
from \citet{DobbinsEtAl2023}. For rational universality we use the
conjunction-only arithmetic ideas in \citet{AbrahamsenMiltzow2019}, with
explicit range checks in Appendix~\ref{app:arithmetic}. The broader
rational-equivalence statement in Theorem~1 of that preprint includes
arbitrary compact semialgebraic sets. Definition~4 of that preprint uses
the same rational-coordinate notion as here: the forward and inverse maps
must be defined on their respective sets. Proposition~\ref{prop:basic-invariance}
and the three-quadrant example rule out that broader scope under this
definition. We therefore prove the required
basic-closed version directly, without its Boolean preprocessing.
Semialgebraic triangulation \citep{OhmotoShiota2017} gives the separate
topological conclusion. The new application is the electrical realization
under all the stated restrictions, not universality of bounded arithmetic
or triangulation itself.

Doubly exponential separation is also known for general polynomial systems:
\citet{JeronimoEtAl2013} give both a polynomial-minimum bound and a
repeated-power example showing the need for that scale. Our contribution
is an explicit infeasible electrical family using only fixed data and
bounded degree, together with a quantitative bound for the actual
arithmetic gadgets. The general issues surrounding short rational points
and approximate certificates are studied by
\citet{BienstockDelPiaHildebrand2023}. Our gap-promise certificate follows
from a direct Lipschitz and rounding argument for the bounded quadratic
network model; we do not present that general approximation principle as
new.

\citet[Theorem 7 and Corollary 8 of the arXiv version]{BienstockMunoz2018}
give LP approximation schemes for polynomial network optimization under
bounded treewidth and bounded numbers of local variables and constraints,
including AC optimal power flow, with scaled feasibility and optimality
tolerances. Those guarantees do not decide exact feasibility.
Our promise result likewise does not resolve the unpromised
approximate-membership question posed for the lossless fixed-magnitude
model in \citet[Section 1.3]{BienstockVermaPreprint2019}.
The quantitative reactive estimate is an explicit consequence of the
energy method, included to delimit the numerical use of the exact AC
transfer. Appendix~\ref{app:verification} records reproducible
exact-arithmetic checks; these support, and do not replace, the proofs.
